\section{Further example economies}
\label{app:further-examples}

This section holds example economies that the main text does not need: the roster of
example economies, the payoff bookkeeping of the mechanism-active $P_3$ economy, the
non-compliance core outside the class, interior compliance at a positive sanction, and
national welfare at the dual-basin economy. The auxiliary polynomial economy
of the roster --- a deliberately simple economy used to test the machinery of
Bellman-invariant sets of value functions and of the joint operator (non-emptiness and invariance of the sets,
dynamic argmax uniqueness, a decoupled Schur contraction) --- is \emph{not} an
example economy at which \Cref{ass:sender,ass:target,ass:concavity} hold jointly: its
sender payoff is linear in the sanction and fails \Cref{ass:concavity}, and
no headline theorem is applied to it. Logical non-vacuity of the full
headline package is supplied by the complete certificate economies of the
roster, including the Eaton--Engers and interest-group constructions. Its
parameter vector and verified invariant-set checks are not reproduced here.

\subsection{Roster of example economies}

\Cref{tab:witness-roster} separates logical non-vacuity,
unconditional applications of the headline results and mechanism displays.
The two conditional existence diagnostics that used to close it --- the
low-$\delta$ purified window and the coupled Schur condition, both at the $P_3$
economy with $\delta=\tfrac1{100}$ --- deliver no equilibrium and carry no
result of this paper; their rows are kept in the notes and their
declarations are indexed in \Cref{app:formalization}.

\begingroup
\captionof{table}{Roster of example economies. The reported fractions
evaluate certified conditions at illustrative economies and are not
calibrated empirical magnitudes.}
\label{tab:witness-roster}
\footnotesize
\setlength{\LTleft}{0pt}
\setlength{\LTright}{0pt}
\begin{longtable}{@{}
  >{\raggedright\arraybackslash}p{0.13\linewidth}
  >{\raggedright\arraybackslash}p{0.22\linewidth}
  >{\raggedright\arraybackslash}p{0.19\linewidth}
  >{\raggedright\arraybackslash}p{0.20\linewidth}
  >{\raggedright\arraybackslash}p{0.12\linewidth}
  @{}}
\toprule
\textbf{Economy} & \textbf{What it certifies} & \textbf{Coupling \& policies} & \textbf{Caveat} & \textbf{Where}\\
\midrule
\endfirsthead
\toprule
\textbf{Economy} & \textbf{What it certifies} & \textbf{Coupling \& policies} & \textbf{Caveat} & \textbf{Where}\\
\midrule
\endhead
\midrule
\multicolumn{5}{r}{\emph{Continued on next page}}\\
\endfoot
\bottomrule
\endlastfoot
Auxiliary polynomial economy
& Invariant set of value functions, unique dynamic argmaxes, and a Schur-contractive joint operator.
& Decoupled; $\alpha\equiv1$ and $\sigma_{\rm dyn}\equiv1$ on the invariant set.
& Linear sender payoff violates A3; no headline-theorem application.
& \Cref{app:witness}\\
\midrule
Eaton--Engers certificate
& Two-sided mechanism-active full-package application: both values are strictly increasing, and continuation reverses each player's named stage ranking over sanctions.
& Full local invariant-set certificate; unconditional persistence, ordering and shock conclusions.
& Decoupled cross-value rows; affine transition; corner policies and corner-pinned realized sanctions.
& \Cref{sec:existence:nonvacuity}\\
\midrule
Quadratic bump
& Same certificate at a non-affine transition.
& Decoupled; corner policies.
& Realized sanctions corner-pinned.
& \Cref{sec:results}\\
\midrule
Integrated nonlinear interior
& One exact nonlinear equilibrium with positive discount factors, interior target and sender actions, strict target sanction preference, and the full persistence floor.
& target policy $\alpha\equiv1/2$; sender sanction rule $(x+1)/2$; quadratic-bump transition.
& Both equilibrium values are zero, so neither forward-looking incentive changes behavior; logical integration, not a substantive calibration.
& \Cref{sec:existence:nonvacuity}\\
\midrule
Integrated nonlinear target-active
& The same nonlinear interior geometry with $V_T(x)=x$; a stage-profitable compliance deviation is strictly rejected by the full target Bellman objective, while the complete persistence floor still holds.
& target policy $\alpha\equiv1/2$; sender sanction rule $(x+1)/2$; quadratic-bump transition.
& $V_S=0$ and cross-value coupling remains absent; only the target's forward-looking incentive changes behavior in this construction.
& \Cref{sec:existence:nonvacuity}\\
\midrule
Low-basin (dovish)
& Mirror low-basin theorem non-vacuous; both invariant directional basins exhibited.
& Decoupled; non-affine transition.
& No interior policy display.
& \Cref{sec:mechanism:bistability}\\
\midrule
Non-monotone (B2)
& Mechanism-active full-package application: positive discounting, nonconstant values, exact nonzero foresight premiums, an interior realized sender sanction, and the full persistence/minimum-drift package; single crossing fails but interval dominance holds.
& At the pivot, $(a^E,s^E)=(0,7/8)$ while the myopic equilibrium has $(a^M,s^M)=(1,5/14)$.
& Cross-value rows are decoupled and the transition is affine; no interior target policy or nonlinear transition in this application.
& \Cref{sec:existence:nonvacuity,rem:ido}\\
\midrule
Interest-group microfounded
& Stage assumptions derived from interest-group primitives.
& Decoupled; corner policies.
& Region~(R); strictly interior basin pivot.
& \Cref{app:fourgroup}\\
\midrule
Derived-basin example
& Inhabits the derived-basin route via its constructor.
& Decoupled; target stage payoff flips at interior pivot $2/5$.
& Kinked sender payoff essential (smooth impossible).
& \Cref{sec:flip}\\
\midrule
Mechanism-active $P_3$
& At $h=1$, an explicit unconditional equilibrium has nonconstant values, a monotone sanction rule, and a strict foresight premium at the top state.
& State-varying leader and sender policies; affine transition; purified box kernel.
& Every $h=1$ equilibrium violates the top-state dynamic floor, so the specialization cannot furnish the three-goals/pathwise-ordering package.
& \Cref{sec:existence:nonvacuity}\\
\midrule
Dual-basin selection
& Both invariant directional basins certified at one economy and the same equilibrium; no limit point is named.
& target row exactly decoupled at the entrenchment-condition boundary; leader policy derived at the corner.
& sender argmax interior only between basins.
& \Cref{sec:mechanism:bistability}\\
\midrule
Quartic economy in the class
& The class inhabited outside the quadratic--linear family: the four cases at $\beta_T=1/2$, the relief equilibrium at $\beta_T=1/30$.
& sender payoff quartic in the sanction, target payoff quadratic in each coordinate; non-affine transition; myopic best response defined as an argmax.
& An existence certificate; no genericity or magnitude claim.
& \Cref{app:witness:quartic}\\
\end{longtable}
\endgroup

\subsection[The P3 example: payoff and cross-partial bookkeeping]
{The $P_3$ example: payoff and cross-partial bookkeeping}

\[
  U_S=a-\tfrac35 s^2-s\big(\tfrac1{10}+\tfrac12 a-\tfrac{27}{20}x\big)+\tfrac25 x,
  \qquad
  U_T=-a^2+2xa-a-s\big(\tfrac12-\tfrac12 a-x\big),
\]
keeping $\delta=\tfrac15,\ \beta_S=\tfrac12,\ \beta_T=\tfrac9{10}$ and the affine
transition. The stronger curvature $-\tfrac35 s^2$ makes the sender's myopic best response
$\sigM=\mathrm{clip}\big(\tfrac98 x-\tfrac1{12}-\tfrac5{12}a\big)$ \emph{interior}
on the lower basin; the target's state--compliance term $+2xa$ makes compliance interior and
reluctant; the sender's stake $+\tfrac25 x$ keeps $\Vs$ state-dependent independently of the
sanction. None of the four changes a structural cross-partial --- the
state--sanction cross-partial signs $\partial^2U_S/\partial s\partial x=\tfrac{27}{20}>0$,
$\partial^2U_T/\partial s\partial x=1>0$, the sign of the compliance--sanction cross term
$\partial^2U_S/\partial s\partial a=-\tfrac12\le0$, and the $s$-concavity all
carry over, and the conditional three-goals
conclusion carries over verbatim.


\paragraph{What the non-affine examples certify, and the transition-free
ordering goal.}
\label{app:witness:grading}
The certified corner recursion is $x_{t+1}=\F(x_t,1)$, so only the corner
section $y\mapsto\F(y,1)$ ever executes along the certified path. The
quadratic-bump transition's bump carries a factor $(1-s)$ that vanishes
there (\path{quadBumpF_corner_isAffine}): that example exercises the
incentive layer under a non-affine transition while its realized path is
the affine one. The low-basin mirror and the wobble economy bend the corner
section itself (\path{quadBumpLowF_corner_not_affine},
\path{wobbleEconF10_corner_not_affine}). All three transitions are refuted
against the affine law at \emph{every} weight, not only at the economy's
own rate (\path{quadBumpF_ne_affineF}, \path{quadBumpLowF_ne_affineF},
\path{ratchetF_ne_affineF}); the non-affine member of the family under the lifting condition is
refuted the same way (\path{quadDiagF_ne_affineF}). Separately, the
certificate route's ordering goal --- the statewise ordering of the
forward-looking sanction above the myopic one --- names no
transition and never needed the affine law: \path{goal2_dominance_generalF}
reaches it with the transition entering only as a lower bound
$d_{\mathrm{lo}}$ on its sanction slope on the target leg and a range-restricted upper bound
$d$ on that slope on the sender leg, at the cost of one sign hypothesis $0\le g_T$
that the affine route does not need because there the transition step is an
equality. The two rates cannot be collapsed into one: a two-sided bound on the sanction slope at a
single rate, with the diagonal condition, pins $\F(x,s)=(1-d)x+ds$. The
affine law is the instance $d_{\mathrm{lo}}=d=\delta$ of
\path{goal2_dominance_generalF}, applied directly at the affine examples (the
separate affine wrapper was removed as redundant in renovation round BG-2).

\paragraph{Open-set robustness at the affine member of the family under the lifting condition.}
\label{app:witness:openbox}
The quadratic--linear family under the lifting condition is inhabited on a box of coefficients
(a statement in the family's coefficients, not in the class's rates; doc 104,
sign-off M6). Let $B$
be the set of coefficient vectors $(q,u_0,k,r,w,\ell,c,n,v)$ with each
coordinate within $1/20$ of the affine member's $(3,1,3,1/2,1,1,4,0,30)$,
except that the target's state--sanction cross term keeps its sign domain, $n\in[0,1/20]$: $B$
is a neighborhood of the member relative to $n\ge0$, not a Euclidean ball
around it, and it contains the full-dimensional open sub-box with $n>0$.
The transition is the affine law at rate $1/2$, the discounts are $1/2$ and
$\sigma_V=1$ throughout, so this is a statement about payoff coefficients,
not an open-set statement in function space. Every point of $B$ is an
economy in the class, and interiority, the lifting condition, the one-period ranking and the strong
entrenchment condition hold there as well: at the worst corners, interiority reads
$u_0+r+(r+w)/3\le2.27<2.95\le q$, the lifting condition the same against $k$, the
one-period ranking $\ell q\le3.21<3.70\le u_0(c-n)$, the strong entrenchment condition
$c\le4.05<9.98\le v/3$, the concavity condition $1/8<q$, and the Riccati closure
$1/8+(r+1/8)^2/(q-1/8)\le1$. On all of $B$ the increasing-differences condition is vacuous
(the transition is flat), both forward-looking incentives change behavior at
the composed equilibrium, in the target-only case the state converges to $u_0/(q-r)$,
and $L_{\mathrm{loop}}=1/2+2(r+1/8)/(q-1/8)<1$, so the closed loop is a
contraction with one globally stable steady state. The sanction-section
modulus is zero, so every point of $B$ is also an economy in the enriched
subclass, on which $M_{\mathrm{tot}}=q$ and the forward-looking relief condition \textup{(RL)} follows from
$2\ell\max\{q,k\}\le1281/200<2301/200\le k(c-n)$ at the worst corners;
hence in the sender-only case the state converges to zero and all four cases of
\Cref{thm:four-cells} are ordered at every point of $B$.

\subsection{The non-compliance core outside the class}
\label{app:witness:outside-chord}

\begin{proposition}[The non-compliance core outside the class]
\label{prop:core-outside-chord}
There is an economy whose target payoff is that of no economy in the class, with an
equilibrium in $\mathcal E$ that has full non-compliance at every state and a nondecreasing
sanction rule, at which every equilibrium in $\mathcal E$ with a nondecreasing sanction
rule has full non-compliance at every state.
\end{proposition}

The economy has the transition $\F(x,s)=(x+s)/2$, both discount factors $\tfrac12$, the
quadratic--linear sender payoff $\Us=x-\tfrac32s^2+(1-3a)s$ without a state--sanction cross term, and the
target payoff $\Ut=3x-s-C(a)$ with compliance cost $C(a)=a+10a(1-a)$. The cost is at
least $a$ but not monotone, $C(3/4)=21/8>C(1)=1$, so the class's requirement that
every compliance increment cost at least $\ell_{\mathrm{lo}}>0$ fails and no economy
in the class has this target payoff (\path{oc_not_chord}). The equilibrium has
$\Vs=4x/3+16/27$, $\Vt=4x$, full non-compliance and the response $\max\{0,4/9-a\}$; the
target's objective is $4x-C(a)$ whatever the sanction, so full non-compliance is its unique
optimum (\path{ocEquilibrium}, \path{oc_argmax_unique}). The conditions of the non-compliance
core (\Cref{app:proofs:threshold}) hold with $g_{\mathrm{lo}}=d_{\mathrm{lo}}=\tfrac12$,
$v=3$ and $c=\ell=1$, and the strong entrenchment condition binds:
$\beta_Td_{\mathrm{lo}}v/(1-\beta_Tg_{\mathrm{lo}})=1=c$ (\path{oc_core_policy_zero},
\path{defiance_core_outside_chord_class}). The economy is an existence certificate: the
core's extra generality over the class is not empty, and nothing is said about how
common such economies are. The transition is affine in this example; the core it
instantiates is general in $\F$, and the generality exhibited is in the target's payoff.

\subsection{Interior compliance at a positive sanction}
\label{app:witness:interior}

\begin{proposition}[A strict interior optimum in the class]
\label{prop:interior-witness}
There is an economy in the class satisfying interiority, the
myopic relief condition, the lifting condition and compliance neutrality, with an
equilibrium in $\mathcal E$ in which the target's compliance $1/14$ is its unique optimum at every
state, its objective falls below that optimum by at least $(a-1/14)^2/2$ at every
compliance $a$, and the sanction is $47/126>0$ at every state.
\end{proposition}

The economy is the affine member of the family under the lifting condition of \Cref{app:witness:openbox} with $r$,
$c$, $v$ and $\beta_T$ set to $0$, $5/2$, $10$ and $1/4$, and the target's payoff
replaced by
$\Ut=10x-\tfrac52s-a-\tfrac12a^2$, whose compliance cost is quadratic. Its compliance
rates lie between $\ell_{\mathrm{lo}}=1$ and $\ell_{\mathrm{hi}}=2$, so the economy is in
the class but outside the family, whose cost is linear (\path{qcChord},
\path{qc_interior}, \path{qc_stageRelief}, \path{qc_placable}, \path{qc_neutral}). The
equilibrium has $\alpha\equiv1/14$, $\Vs=4x/3+2209/5292$ and $\Vt=80x/7-557/882$
(\path{qcEquilibrium}, \path{qc_rule}). The growth bound is \path{qc_growth}, so no
first-order condition or grid is used (\path{qc_argmax_unique},
\path{strict_interior_equilibrium_in_chord_class}). On the family the cost is linear
and interior compliance is an equilibrium only through indifference; the curvature of
the cost is what makes the interior optimum strict. The economy is an existence
certificate, with an affine transition, and says nothing about how common such economies
are.

\subsection{National welfare at the dual-basin economy}
\label{app:entrenchment-welfare}

At the dual-basin economy each government's payoff is its nation's welfare plus a
political part that grows with the political state. Define the sender's and the
target's national welfare by
\[
W_S(a,s)=a-\tfrac12s^2-\tfrac12s-2as,\qquad
W_T(a,s)=-\tfrac15a-\tfrac{9}{10}s ,
\]
so that $\Us(x,a,s)=W_S(a,s)+2xs$ and $\Ut(x,a,s)=W_T(a,s)+\tfrac{19}{10}x$. Both
political parts vanish at the state $x=0$ and rise with the state. The statement
was in the main text until 2026-10-15.

\begin{proposition}[Entrenchment against both nations]
\label{prop:entrenchment-welfare}
At the dual-basin economy:
\begin{enumerate}[label=(\roman*)]
\item At the state $x=0$, each government's payoff equals its nation's welfare
at every action. At every state $x\le1/4$ the sender's myopic best response to
every compliance $a\ge0$ is no sanction.
\item There is an equilibrium in $\mathcal E$ such that, from every initial
state, after a temporary shock that holds the sanction at $1$ for at least
$\widehat N_{\F}(19/20)$ periods, at every later date the state lies in
$[19/20,1]$, the target maintains full non-compliance and the sanction is $1$. In every such
period the sender's nation has welfare $-1$ and the target's $-9/10$, against
$0$ for both at full trade, so at every discount factor $\beta\in[0,1)$ their
discounted welfare is $-1/(1-\beta)$ and $-(9/10)/(1-\beta)$. The sender
government's payoff in every such period is $2x_t-1\ge 9/10$, above its payoff
$0$ at full trade. At its own discount factor, the target government's payoff
from the state $x$ at which the shock ends is $10x$, the same as under
permanent relief from that state.
\item Replace the sender's payoff by its nation's welfare $W_S$, with everything
else unchanged. The resulting economy has an equilibrium in $\mathcal E$ at
which the sanction is zero at every state. Every equilibrium in $\mathcal E$ of
that economy whose target value rises with the state at exactly rate $10$ and
is bounded by $19$ in absolute value, the set of target values used to
construct \Cref{prop:two-sided-one-equilibrium}, has zero sanctions at every
state.
\end{enumerate}
\end{proposition}

\begin{proof}[Proof of \Cref{prop:entrenchment-welfare}]
(i) The identities $\Us=W_S+2xs$ and $\Ut=W_T+\tfrac{19}{10}x$ hold at every
state and action, and both political parts are zero at $x=0$ and
non-decreasing in $x$ for $s\ge0$ (\path{dbU_S_split}, \path{dbU_T_split},
\path{db_political_parts_zero_at_full_trade},
\path{db_political_parts_monotone}). The sender's myopic best response is the
clipped vertex $\max\{0,\min\{1,2x-\tfrac12-2a\}\}$, which is $0$ when
$x\le1/4$ and $a\ge0$ (\path{db}$\sigma$\path{_eq_zero_of_le}).

(ii) Take an equilibrium of the dual-basin economy whose values lie in the
value sets of the high-basin certificate, as for
\Cref{prop:two-sided-one-equilibrium}. After the shock, the resumed path stays
in $[19/20,1]$, where the target's policy is full non-compliance and the realized
sanction is $1$. Hence the national welfare in each period is
$W_S(0,1)=-1$ and $W_T(0,1)=-9/10$, while both are $0$ at full trade, and the
geometric sums give $-1/(1-\beta)$ and $-(9/10)/(1-\beta)$. The sender
government's payoff is $\Us(x,0,1)=2x-1\ge9/10$ on the basin, against
$\Us(0,0,0)=0$. Under permanent full sanctions from a state $z$ the state is
$1-(1-z)(9/10)^n$, so at $\beta_T=9/10$ the target government's payoff is
\[
\sum_{n\ge0}\bigl(\tfrac{9}{10}\bigr)^n
\Bigl[\tfrac{19}{10}\bigl(1-(1-z)(\tfrac{9}{10})^n\bigr)-\tfrac{9}{10}\Bigr]
=\tfrac{19}{10}\Bigl(10-(1-z)\tfrac{100}{19}\Bigr)-9=10z .
\]
Under permanent relief from $z$ the state is $z(9/10)^n$, and the payoff is
$\sum_{n\ge0}(9/10)^n\,\tfrac{19}{10}z(9/10)^n=\tfrac{19}{10}z\cdot\tfrac{100}{19}=10z$
(\path{db_entrenchment_loss_after_episode}).

(iii) With $\Us$ replaced by $W_S$, the sender's stage payoff against full
non-compliance, $W_S(0,s)=-\tfrac12s^2-\tfrac12s$, does not depend on the state, is
$0$ at $s=0$ and negative at every $s\in(0,1]$. Against a target policy of full
non-compliance at every state, such a sender's value is identically zero and its
forward-looking best response is $0$ at every state
(\path{dynamicTopRealizedRule_eq_zero_of_stateFree}, a general statement in
\path{Theory/Welfare/RepresentativeSender.lean}). The values $\Vs\equiv0$ and
$\Vt(x)=10x$ with full non-compliance and no sanction form an equilibrium in
$\mathcal E$: $10x=\tfrac{19}{10}x+\tfrac{9}{10}\cdot10\cdot\tfrac{9x}{10}$, and the
target's objective $10x-a/5$ is maximal at $a=0$ (\path{dualRepEquilibriumTop},
\path{dualRep_rule_eq_zero}). For an equilibrium whose target value is
$\Vt(x)=c+10x$, the target's objective is
$\tfrac{19}{10}x-\tfrac15a-\tfrac{9}{10}s+\tfrac{9}{10}\bigl(c+(9x+s)\bigr)
=10x-\tfrac15a+\tfrac{9}{10}c$: a sanction costs the target exactly what the
entrenchment it builds is worth to it, so the objective does not depend on the
sender's choice, the target maintains full non-compliance, and the sender side is as above
(\path{dualRep_rule_eq_zero_of_carrier}).
\end{proof}
